\section{A polynomial odd Laplacian}

\begin{proposition}[A polynomial Batalin--Vilkovisky operator]
\label{prop:cbc26-polynomial-laplacian}
On $K=\mathbb C[x_1,\ldots,x_n]\otimes\Lambda(\eta_1,\ldots,\eta_n)$, with $|x_i|=0$ and $|\eta_i|=-1$, use left odd derivatives and put
\[
 \Delta=\sum_i\partial_{x_i}\partial_{\eta_i}.
\]
This operator has degree one, order at most two, and square zero.
For homogeneous $f,g$, define
\[
 \{f,g\}=(-1)^{|f|}\bigl(\Delta(fg)-(\Delta f)g-(-1)^{|f|}f\Delta g\bigr).
\]
This odd bracket has $\{x_i,\eta_j\}=\delta_{ij}$.
For every polynomial $S(x)$, one has $\Delta S=\{S,S\}=0$, and $D=\{S,-\}$ satisfies $D\Delta+\Delta D=0$.
Thus $S$ solves the polynomial quantum master equation $\tfrac12\{S,S\}+\hbar\Delta S=0$.
\end{proposition}
\begin{proof}
The even derivatives commute and distinct odd derivatives anticommute; equal odd derivatives square to zero.
Pairing indices in $\Delta^2$ proves its vanishing.
Each summand is a composition of two derivations and hence has order at most two.
Expanding $\Delta(fg)$ gives the degree-one derived bracket; evaluation on generators yields the stated brackets, which determine its biderivation extension.
Both $\Delta S$ and $\{S,S\}$ vanish since $S$ contains no $\eta$.
In $D\Delta+\Delta D$, the terms with a derivative left on the input cancel by odd anticommutation, while the remaining Hessian coefficients $\partial_i\partial_jS$ are symmetric and multiply the antisymmetric derivatives $\partial_{\eta_i}\partial_{\eta_j}$.
Their sum is zero.
\end{proof}

\begin{remark}[A chiral Batalin--Vilkovisky comparison]\label{thm:bv-structure-bar}
The polynomial operator above is defined on its displayed finite generator algebra.
A residue operator on a chiral configuration complex requires a domain, a cyclic pairing, its degree and second-order identity, and compatibility with every collision map.
The Arnold relations alone do not specify these data or an interacting master action.
An identification with the polynomial model requires the actual coefficient and operator maps.
\end{remark}
